%mac
\subsubsection{TỐC ĐỘ}
\paragraph{Tốc độ trung bình}
\textit{Tốc độ trung bình} là đại lượng vật lí đặc trưng cho sự nhanh, chậm của chuyển động trong một khoảng thời gian.
\begin{align*}
\boder{v_{tb}=\dfrac{s}{t}}\quad \ct{với} \begin{cases} s: \ct{quãng đường đi được (m; km)}\\
t: \ct{thời gian chuyển động (s; h)}\end{cases}
\end{align*}
%\newghichu{Nếu quãng đường đi được tại thời điểm $t_1$ là $s_1$, quãng đường đi được tại thời điểm $t_2$ là $s_2$ thì:\\
%	\begin{align*}v_{tb}=\dfrac{\Delta s}{\Delta t}=\dfrac{s_2-s_1}{t_2-t_1}\end{align*}}
\paragraph{Chuyển động thẳng đều}
\textit{Chuyển động thẳng đều} là chuyển động có quỹ đạo là đường thẳng và có tốc độ trung bình như nhau trên mọi quãng đường.
\paragraph{Quãng đường đi được trong chuyển động thẳng đều}
\begin{align*}
\boder{s=v.t}
\end{align*}
\Note{
Trong chuyển động thẳng đều, tốc độ là không đổi, quãng đường đi tỉ lệ thuận với thời gian chuyển động.
}
\subsubsection{VẬN TỐC}
\setcounter{paragraph}{0}
\paragraph{Vận tốc trung bình}
\begin{align*}
\boder{\vec v=\dfrac{\vec d}{t}}
\end{align*}
\begin{itemize}
\item Vì độ dịch chuyển là một đại lượng vectơ nên vận tốc cũng là một đại lượng vectơ. 
\item Vectơ vận tốc có: 
\begin{itemize}
\item Gốc nằm trên vật chuyển động.
\item Hướng là hướng của độ dịch chuyển.
\item Độ dài tỉ lệ với độ lớn của vận tốc.
\end{itemize}
\end{itemize}
\paragraph{Vận tốc tức thời} 
\begin{itemize}
\item Vận tốc tức thời là vận tốc tại một thời điểm xác định, được kí hiệu $\vec v_t$:
\begin{align*}
\vec v_t=\dfrac{\Delta \vec d}{\Delta t}\ \text{với}\ \Delta t \ \text{rất nhỏ}
\end{align*}
\end{itemize}


